\begin{textsample}{POS Dim 1 – vem\_ed\_13 – Score 13.00 – vem\_ed\_13/t057.txt}  \label{ex:f1_pos_124}
Coragem para decidir Ricardo Nóbrega \textbf{é} coordenador do curso de Comércio Exterior da Fatec Indaiatuba. \textbf{Começou} o primeiro PCI em 2019, com a DePaul University (EUA).
Atualmente, além da DePaul, a Fatec Indaiatuba tem PCIs com as seguintes instituições: Universidade Católica de Hong Kong, Instituto Politécnico de Viseu (Portugal), Durban University of Technology (África do Sul), Eastern Oregon University (EUA), Universidad Politécnica San Luís de Potosí (México) e Amsterdam University of Applied Sciences (Holanda).
A \textbf{proposta} de Nóbrega \textbf{é} distribuir os PCIs ao \textbf{longo} dos semestres, em uma sequência ordenada, para eliminar a barreira do \textbf{idioma}.
Os alunos iniciam sua experiência com Portugal.
A \textbf{seguir}, participam de PCIs em língua espanhola e, por fim, inglesa.
Em junho de 2021, Nóbrega apresentou trabalho sobre os PCIs na unidade durante a conferência"Going Global".
A \textbf{seguir}, trechos de seu depoimento.
Os PCIs tiram professores e alunos da zona de conforto.
Eles vão muito além do \textbf{aprendizado} de \textbf{idiomas}, têm um potencial enorme, muito maior do \textbf{que} a gente possa imaginar, desenvolvendo nos participantes a \textbf{capacidade} de adaptação, comunicação, solução de problemas e negociação. \textbf{São} projetos com uma complexidade muito grande. \textbf{É} muito diferente trabalhar com os Estados Unidos ou com a África do Sul: o perfil dos alunos e os \textbf{contextos}, as \textbf{realidades}, \textbf{são} completamente distintos.
Os PCIs possuem uma dinâmica interessante, com \textbf{mudanças} assim \textbf{que} \textbf{começam}, o \textbf{que} exige \textbf{flexibilidade} para as adaptações necessárias.
Tanto \textbf{é} \textbf{que} estou substituindo a expressão"planejamento"por"condições iniciais".
Mais vale a \textbf{capacidade} de se adaptar e ter a coragem de tomar decisões, para \textbf{que} o projeto não pare.
Além disso, organização e gestão do tempo \textbf{são} grandes desafios [ \textbf{são} os mais citados por professores e alunos participantes dos PCIs, nas pesquisas de percepção realizadas desde 2020 ].
Para organizar e distribuir a divisão das \textbf{tarefas}, \textbf{cada} grupo precisa \textbf{identificar} os diversos talentos: quem fala melhor, quem tem mais disponibilidade, os \textbf{que} \textbf{são} bons em pesquisa, em texto... Às vezes os alunos se queixam de falhas na comunicação, mas nesses projetos gosto muito de pensar no binômio comunicação versus postura.
Alguns alunos precisam trabalhar a maturidade, o medo de se expor.
Precisamos estimular a confiança dos estudantes.
Para reservar um espaço dedicado aos PCIs, estamos \textbf{criando} a"Global Room"na Fatec Indaiatuba, com dois televisores e computadores \textbf{conectados} à internet. (Depoimento gravado em reunião com a equipe PCI / Cesu, organizada na Plataforma Teams em 3 de maio de 2022).

% matched lemmas: aprendizado, cada, capacidade, começar, conectar, contexto, criar, flexibilidade, identificar, idioma, longo, mudança, proposta, que, realidade, seguir, ser, tarefa
\end{textsample}
